Two things stand out.

\textbf{More words took more sets, but far fewer than in proportion.} From 16 to 1,024 animals
there are 64 times as many facts. Taking the middle one of the three figures in each entry of
table \ref{tab:effort1}, width 8 needed 8 times as many sets (512 against 64), and widths 16
and 32 needed about 6 times as many.

\textbf{The smaller the model, the more sets it needed, and the less dependable it was.} On
1,024 animals, width 8 needed four times as many sets as width 32 (512 against 128, the middle
figures), and its
three repeats differed widely: 256, 512 and 1,639 sets. The slowest only just made it: it
passed for the second time at the grading made when its 30 minutes ran out. For its first 250 sets its score stayed near zero. Later it stood still
again, at a score of 25 from set 512 to set 724, before climbing to 52 and then to 100. (A
score of 25 is what a model gets if it knows one of the four kinds of fact for every animal, so
it seems to have picked up one kind of fact at a time.) This is a warning for everything that
follows: a model can stand still for a long time and then move again.
